\begingroup
\setlength\LTleft{0\linewidth}
\setlength\LTright{0\linewidth}
\fontsize{6.0pt}{7.2pt}\selectfont

\setlength{\LTpost}{0mm}
\begin{longtable}{@{\extracolsep{\fill}}lrrrrrrr}

\toprule
 &  &  & \multicolumn{2}{c}{\textbf{1ª metade}} & \multicolumn{2}{c}{\textbf{Pós-promoção}} &  \\ 
\cmidrule(lr){4-5} \cmidrule(lr){6-7}
\textbf{Variante} & \textbf{Itens} & \textbf{Retorno composto} & \textbf{Coef.} & \textbf{p} & \textbf{Coef.} & \textbf{p} & \textbf{ROI da carteira} (3) \\ 
\midrule\addlinespace[2.5pt]
Base & 10.515 & 0,17\% & -0,59\% & 0,004 & 0,37\% & 0,035 & 43,61\% \\
Cobertura mínima de 30\% (1) & 10.515 & 0,14\% & -0,56\% & 0,006 & 0,40\% & 0,022 & 52,66\% \\
Cobertura mínima de 70\% (1) & 10.496 & 0,18\% & -0,59\% & 0,002 & 0,35\% & 0,061 & 44,03\% \\
Mínimo de 100 dias (2) & 8.689 & 0,16\% & -0,59\% & 0,004 & 0,37\% & 0,035 & 43,61\% \\
Mínimo de 200 dias (2) & 6.588 & 0,15\% & -0,56\% & 0,007 & 0,37\% & 0,029 & 53,66\% \\
Mínimo de 350 dias (2) & 4.719 & 0,11\% & -0,56\% & 0,006 & 0,38\% & 0,033 & 60,06\% \\
Mínimo de 500 dias (2) & 366 & 0,03\% & -0,68\% & 0,002 & 0,46\% & 0,011 & Inviável \\
Intervalo máximo de 3 dias & 10.515 & 0,12\% & -0,60\% & 0,006 & 0,43\% & 0,028 & 43,21\% \\
Teto do retorno em 100\% & 10.515 & 0,06\% & -0,62\% & 0,001 & 0,38\% & 0,023 & 42,55\% \\
Teto do retorno em 200\% & 10.515 & 0,13\% & -0,60\% & 0,002 & 0,37\% & 0,031 & 43,61\% \\
Teto do retorno em 1.000\% & 10.515 & 0,18\% & -0,58\% & 0,004 & 0,37\% & 0,033 & 43,61\% \\
Deslocamento do calendário (dias): −1 & 10.515 & 0,17\% & -0,56\% & 0,004 & 0,35\% & 0,056 & 43,61\% \\
Deslocamento do calendário (dias): +1 & 10.515 & 0,17\% & -0,62\% & 0,002 & 0,37\% & 0,019 & 43,61\% \\
Janela de preços (observações): 1 & 10.515 & 0,17\% & -0,59\% & 0,004 & 0,37\% & 0,035 & 39,30\% \\
\bottomrule
\end{longtable}
\begin{minipage}{\linewidth}
Fonte: Elaboração própria.\\
Coeficientes do retorno diário do mercado com ponderação igual, como na tabela de fases.\\
(1) Cobertura mínima: parcela mínima da mediana diária de preços positivos que o dia precisa atingir.\\
(2) Mínimo de dias com preço no período de treinamento para o item entrar na amostra.\\
(3) ROI líquido com o tamanho comum de 50 itens, em ordem decrescente do ROI no histórico de seleção. Inviável quando menos de 50 itens atendem aos critérios.\\
\end{minipage}
\endgroup
